\documentclass[10pt,a4paper]{amsart}
\usepackage[margin=19mm]{geometry}
\usepackage{amsmath,amssymb,mathtools}
\usepackage[scr=rsfs,cal=euler]{mathalfa}
\usepackage{xfrac}
\usepackage[unicode,hidelinks,bookmarks=false]{hyperref}
\newcommand{\ie}{i.e.\ }
\newcommand{\cf}{cf.\ }
\newcommand{\resp}{respectively\ }
\newcommand{\Subsection}{\subsection}
\providecommand{\nobreakdash}{\nobreak\hbox{-}}
\setlength{\emergencystretch}{2em}
\multlinegap0pt
\usepackage{enumitem}
\DeclareMathOperator{\csch}{csch}
\DeclareMathOperator{\sech}{sech}
\newtheorem{theorem}{Theorem}[section]
\newtheorem{lemma}[theorem]{Lemma}
\newtheorem{proposition}[theorem]{Proposition}
\newtheorem{corollary}[theorem]{Corollary}
\newtheorem{conjecture}[theorem]{Conjecture}
\newtheorem{definition}[theorem]{Definition}
\newtheorem{remark}[theorem]{Remark}
\begin{document}
\title[A Heat Kernel Expectation Approach to Boundary-Corrected Li-]{A Heat Kernel Expectation Approach to Boundary-Corrected Li--Yau Estimates for the Dirichlet Heat Equation}
\author{Li-Chang Hung}
\date{}
\maketitle
\noindent\emph{Source and licence.} A Heat Kernel Expectation Approach to Boundary-Corrected Li--Yau Estimates for the Dirichlet Heat Equation. Version arXiv:2608.12376v1 (CC BY 4.0). This material is distributed under the Creative Commons Attribution 4.0 International licence, http://creativecommons.org/licenses/by/4.0/, as recorded in the arXiv OAI licence metadata for this exact version and shown on the abstract page https://arxiv.org/abs/2608.12376v1. Full rights notice: licenses/arxiv-2608.12376-CC-BY-4.0.txt.\par\smallskip
\noindent\textit{Editorial scope.} Counted unit: the original \texttt{proposition} environment \texttt{-} at source lines 1754--1775 together with its complete proof at lines 1778--1832; the delivered document keeps source lines 728--1980 so that every referenced label and every citation resolves inside it. Native editable source is the paper's own concatenated TeX; the AMS wrapper is editorial formatting only. No publisher class, upstream script or unrelated artwork is redistributed.
\section{Proof of the Main Theorem}
%%%%%%%%%%%%%%%%%%%%%%%%%%%%%%%%%%%%%%%%%%%%%%%%


We compute each component of

\[
\Delta\log w
=
\sum_{i=1}^{n}
\left(
\frac{w_{x_ix_i}}w
-
\frac{w_{x_i}^2}{w^2}
\right).
\]


For each coordinate,

\[
\frac{w_{x_ix_i}}w
=
\int
\frac{\partial_{x_ix_i}K_D}
{K_D}
d\mu_{x,t}(y).
\]



%%%%%%%%%%%%%%%%%%%%%%%%%%%%%%%%%%%%%%%%%%%%%%%%
\subsection{Tangential contribution}
%%%%%%%%%%%%%%%%%%%%%%%%%%%%%%%%%%%%%%%%%%%%%%%%


For $i<n$,

\[
\partial_{x_ix_i}K_D
=
\left(
\frac{(x_i-y_i)^2}{4t^2}
-\frac1{2t}
\right)K_D .
\]

Hence,

\[
\frac{w_{x_ix_i}}w
=
\int
\left(
\frac{(x_i-y_i)^2}{4t^2}
-\frac1{2t}
\right)
d\mu_{x,t}.
\]


Using Jensen,

\[
\frac{w_{x_ix_i}}w
-
\frac{w_{x_i}^2}{w^2}
\geq
-\frac1{2t}.
\]


Summing over

\[
i=1,\ldots,n-1,
\]

gives

\[
\sum_{i=1}^{n-1}
\left(
\frac{w_{x_ix_i}}w
-
\frac{w_{x_i}^2}{w^2}
\right)
\geq
-\frac{n-1}{2t}.
\]



%%%%%%%%%%%%%%%%%%%%%%%%%%%%%%%%%%%%%%%%%%%%%%%%
\subsection{Normal contribution}
%%%%%%%%%%%%%%%%%%%%%%%%%%%%%%%%%%%%%%%%%%%%%%%%


For the normal variable,

\[
\partial_{x_nx_n}K_D
=
\frac1{4t^2}
\left(
(x_n-y_n)^2-\!2t
\right)
e^{-\frac{(x_n-y_n)^2}{4t}}
\]

\[
-
\frac1{4t^2}
\left(
(x_n+y_n)^2-\!2t
\right)
e^{-\frac{(x_n+y_n)^2}{4t}} .
\]


After dividing by $K_D$ and simplifying, we obtain

\[
\frac{\partial_{x_nx_n}K_D}
{K_D}
=
\frac1{4t^2}
\left(
-2t+x_n^2+y_n^2
-
2x_ny_n
\coth
\left(
\frac{x_ny_n}{2t}
\right)
\right).
\]


Therefore,

\[
\frac{w_{x_nx_n}}w
=
\int
\frac1{4t^2}
\left(
-2t+x_n^2+y_n^2
-
2x_ny_n
\coth
\frac{x_ny_n}{2t}
\right)
d\mu .
\]


On the other hand,

\[
\frac{w_{x_n}}w
=
\int
\left(
\frac{y_n}{2t}
\coth
\frac{x_ny_n}{2t}
-
\frac{x_n}{2t}
\right)
d\mu .
\]


Applying Jensen,

\[
\left(\frac{w_{x_n}}w\right)^2
\leq
\int
\left(
\frac{y_n}{2t}
\coth
\frac{x_ny_n}{2t}
-
\frac{x_n}{2t}
\right)^2
d\mu .
\]


Hence,

\[
\begin{aligned}
&
\frac{w_{x_nx_n}}w
-
\frac{w_{x_n}^2}{w^2}
\\
&\geq
\int
\frac1{4t^2}
\left(
-2t
+
y_n^2
-
y_n^2
\coth^2
\frac{x_ny_n}{2t}
\right)
d\mu .
\end{aligned}
\]


Using

\[
1-\coth^2 z
=
-\csch^2 z ,
\]

we get

\[
\begin{aligned}
&
\frac{w_{x_nx_n}}w
-
\frac{w_{x_n}^2}{w^2}
\\
&\geq
\int
\left(
-\frac1{2t}
-
\frac{y_n^2}{4t^2}
\csch^2
\frac{x_ny_n}{2t}
\right)d\mu .
\end{aligned}
\]
\iffalse
\[
\geq
\int
\left(
-\frac1{2t}
-
\frac{y_n^2}{4t^2}
\csch^2
\frac{x_ny_n}{2t}
\right)d\mu .
\]
\fi

Now

\[
\frac{y_n^2}{4t^2}
\csch^2
\left(
\frac{x_ny_n}{2t}
\right)
=
\frac1{x_n^2}
\left(
\frac{x_ny_n}{2t}
\csch
\frac{x_ny_n}{2t}
\right)^2 .
\]


Since

\[
0<z\csch z<1,
\]

we obtain

\[
\frac{y_n^2}{4t^2}
\csch^2
\left(
\frac{x_ny_n}{2t}
\right)
\leq
\frac1{x_n^2}.
\]


Therefore,

\[
\frac{w_{x_nx_n}}w
-
\frac{w_{x_n}^2}{w^2}
\geq
-\frac1{2t}
-\frac1{x_n^2}.
\]


Combining tangential and normal estimates,

\[
\Delta\log w
\geq
-\frac{n-1}{2t}
-\frac1{2t}
-\frac1{x_n^2}.
\]

Hence,

\[
\boxed{
\Delta\log w
\geq
-\frac n{2t}
-\frac1{x_n^2}
}.
\]

The proof is complete.

\hfill$\square$


%%%%%%%%%%%%%%%%%%%%%%%%%%%%%%%%%%%%%%%%%%%%%%%%
\section{Boundary-Corrected Harnack Inequality}
%%%%%%%%%%%%%%%%%%%%%%%%%%%%%%%%%%%%%%%%%%%%%%%%

The differential Li--Yau estimate obtained in the previous section
immediately yields an integrated Harnack inequality.

Let

\[
l(x,t)=\log w(x,t).
\]

Since

\[
w_t=\Delta w,
\]

we have

\[
l_t
=
\frac{w_t}{w}
=
\frac{\Delta w}{w}.
\]

On the other hand,

\[
\Delta l
=
\frac{\Delta w}{w}
-
\frac{|\nabla w|^2}{w^2},
\]

and therefore

\[
l_t
=
\Delta l+|\nabla l|^2 .
\]


Combining this identity with the boundary-corrected Li--Yau
estimate

\[
\Delta l
\geq
-\frac n{2t}
-\frac1{x_n^2},
\]

we obtain

\[
l_t
\geq
|\nabla l|^2
-\frac n{2t}
-\frac1{x_n^2}.
\]


The following theorem gives the corresponding integral Harnack
inequality.

%%%%%%%%%%%%%%%%%%%%%%%%%%%%%%%%%%%%%%%%%%%%%%%%

\begin{theorem}[Boundary-corrected Harnack inequality]

Let $w$ be a positive solution of the Dirichlet heat equation

\[
w_t=\Delta w
\]

on the Euclidean half-space

\[
\mathbb R^n_+
=
\{x_n>0\},
\]

and assume that

\[
\Delta\log w(x,t)
\geq
-\frac n{2t}
-\frac1{x_n^2}.
\]

Then for every smooth curve

\[
\gamma(t)=(x(t),t),
\qquad
t\in[t_1,t_2],
\]

with

\[
x(t)\in\mathbb R^n_+,
\]

we have

\[
\begin{aligned}
\log
\frac{w(x(t_2),t_2)}
     {w(x(t_1),t_1)}
\geq&
-\frac14
\int_{t_1}^{t_2}
|x'(t)|^2\,dt
\\
&-\frac n2
\log\frac{t_2}{t_1}
\\
&-
\int_{t_1}^{t_2}
\frac1{x_n(t)^2}\,dt .
\end{aligned}
\]

\end{theorem}


%%%%%%%%%%%%%%%%%%%%%%%%%%%%%%%%%%%%%%%%%%%%%%%%

\begin{proof}

Let

\[
l=\log w .
\]

Along a smooth curve $x(t)$ in the half-space, we have

\[
\frac{d}{dt}l(x(t),t)
=
l_t+\nabla l\cdot x'(t).
\]

Using the differential inequality above,

\[
l_t
\geq
|\nabla l|^2
-\frac n{2t}
-\frac1{x_n^2},
\]

we obtain

\[
\frac{d}{dt}l(x(t),t)
\geq
|\nabla l|^2
+
\nabla l\cdot x'(t)
-
\frac n{2t}
-\frac1{x_n^2}.
\]

The first two terms can be estimated by completing the square:

\[
|\nabla l|^2
+
\nabla l\cdot x'
=
\left|
\nabla l+\frac12x'
\right|^2
-
\frac14|x'|^2 .
\]

Therefore,

\[
|\nabla l|^2
+
\nabla l\cdot x'
\geq
-\frac14|x'|^2 .
\]

Consequently,

\[
\frac{d}{dt}l(x(t),t)
\geq
-\frac14|x'(t)|^2
-\frac n{2t}
-\frac1{x_n(t)^2}.
\]

Integrating from $t_1$ to $t_2$ gives

\[
\begin{aligned}
l(x(t_2),t_2)-l(x(t_1),t_1)
\geq&
-\frac14
\int_{t_1}^{t_2}
|x'(t)|^2dt
\\
&
-\frac n2
\int_{t_1}^{t_2}
\frac1t dt
\\
&
-
\int_{t_1}^{t_2}
\frac1{x_n(t)^2}dt .
\end{aligned}
\]

Since

\[
\int_{t_1}^{t_2}\frac1t dt
=
\log\frac{t_2}{t_1},
\]

and

\[
l=\log w,
\]

we obtain

\[
\begin{aligned}
\log
\frac{w(x(t_2),t_2)}
     {w(x(t_1),t_1)}
\geq&
-\frac14
\int_{t_1}^{t_2}
|x'(t)|^2dt
\\
&
-\frac n2
\log\frac{t_2}{t_1}
\\
&
-
\int_{t_1}^{t_2}
\frac1{x_n(t)^2}dt .
\end{aligned}
\]

This proves the theorem.

\end{proof}


%%%%%%%%%%%%%%%%%%%%%%%%%%%%%%%%%%%%%%%%%%%%%%%%

\begin{remark}

When the boundary correction is removed, namely in the whole
Euclidean space, the above estimate reduces to the classical
Li--Yau Harnack inequality

\[
\log
\frac{u(x_2,t_2)}
     {u(x_1,t_1)}
\geq
-\frac{|x_2-x_1|^2}{4(t_2-t_1)}
-\frac n2
\log\frac{t_2}{t_1}.
\]

The additional term

\[
-\int_{t_1}^{t_2}
\frac1{x_n(t)^2}\,dt
\]

is the explicit contribution of the absorbing boundary.

\end{remark}


%%%%%%%%%%%%%%%%%%%%%%%%%%%%%%%%%%%%%%%%%%%%%%%%

\begin{remark}

The boundary correction term has a geometric interpretation.

Since

\[
x_n=d(x,\partial\mathbb R^n_+),
\]

the Harnack inequality can be written as

\[
-\int_{t_1}^{t_2}
\frac1{
d(x(t),\partial\mathbb R^n_+)^2
}
dt .
\]

Thus the influence of the boundary becomes stronger when the
diffusion path approaches the boundary.

This suggests a natural extension to general convex domains,
where the distance function to the boundary is expected to replace
$x_n$.

\end{remark}


%%%%%%%%%%%%%%%%%%%%%%%%%%%%%%%%%%%%%%%%%%%%%%%%
\section{Optimal Paths and Boundary Geodesics}
%%%%%%%%%%%%%%%%%%%%%%%%%%%%%%%%%%%%%%%%%%%%%%%%


The Harnack inequality obtained in the previous section naturally
leads to a variational problem describing the optimal paths of heat
propagation.

Indeed, the boundary-corrected Harnack inequality gives

\[
\log
\frac{w(x(t_2),t_2)}
     {w(x(t_1),t_1)}
\geq
-\int_{t_1}^{t_2}
\left(
\frac14 |x'(t)|^2
+
\frac1{x_n(t)^2}
\right)dt
-
\frac n2
\log\frac{t_2}{t_1}.
\]


Therefore, the optimal path connecting two points is determined by
minimizing the action functional

\[
\mathcal A[x]
=
\int_{t_1}^{t_2}
\left(
\frac14 |x'(t)|^2
+
\frac1{x_n(t)^2}
\right)dt .
\]


The first term represents the classical kinetic energy of diffusion,
while the second term is a new potential term generated by the
Dirichlet boundary.

This variational problem can be viewed as a boundary-corrected
geodesic problem associated with the Li--Yau estimate.



%%%%%%%%%%%%%%%%%%%%%%%%%%%%%%%%%%%%%%%%%%%%%%%%
\subsection{Euler--Lagrange equation}


Let

\[
L(x,x')
=
\frac14|x'|^2+\frac1{x_n^2}.
\]


The Euler--Lagrange equation is

\[
\frac{d}{dt}
\frac{\partial L}{\partial x_i'}
-
\frac{\partial L}{\partial x_i}
=0.
\]


For the tangential directions

\[
i=1,\ldots,n-1,
\]

we have

\[
\frac{\partial L}{\partial x_i'}
=
\frac12 x_i',
\]

and

\[
\frac{\partial L}{\partial x_i}=0 .
\]

Hence,

\[
\boxed{
x_i''=0,
\qquad i=1,\ldots,n-1 .
}
\]


Therefore, the motion parallel to the boundary is a straight line
with constant velocity.

The nontrivial behavior occurs in the normal direction.



%%%%%%%%%%%%%%%%%%%%%%%%%%%%%%%%%%%%%%%%%%%%%%%%
\subsection{Normal boundary direction}


For the normal component,

\[
x_n=d(x,\partial\mathbb R^n_+),
\]

the reduced Lagrangian is

\[
L_n
=
\frac14(x_n')^2
+
\frac1{x_n^2}.
\]


The Euler--Lagrange equation becomes

\[
\frac d{dt}
\frac{\partial L_n}{\partial x_n'}
-
\frac{\partial L_n}{\partial x_n}
=0 .
\]


Since

\[
\frac{\partial L_n}{\partial x_n'}
=
\frac12x_n',
\]

we have

\[
\frac d{dt}
\frac{\partial L_n}{\partial x_n'}
=
\frac12x_n'' .
\]


Moreover,

\[
\frac{\partial L_n}{\partial x_n}
=
-2x_n^{-3}.
\]


Consequently,

\[
\frac12x_n''
-
(-2x_n^{-3})
=
0 .
\]


Hence the optimal normal paths satisfy

\[
\boxed{
x_n''
+
4x_n^{-3}
=
0 .
}
\]


This equation coincides with the optimal path equation obtained by
Hamilton for the one-dimensional Dirichlet heat equation
\cite{Hamilton2011}.



%%%%%%%%%%%%%%%%%%%%%%%%%%%%%%%%%%%%%%%%%%%%%%%%
\subsection{Energy conservation}


The equation

\[
x_n''
+
4x_n^{-3}
=
0
\]

admits a conserved quantity.

Multiplying by $x_n'$ gives

\[
x_n'x_n''
+
4x_n^{-3}x_n'
=
0 .
\]


Since

\[
x_n'x_n''
=
\frac12
\frac d{dt}(x_n')^2,
\]

and

\[
x_n^{-3}x_n'
=
-\frac12
\frac d{dt}x_n^{-2},
\]

we obtain

\[
\frac d{dt}
\left(
\frac12(x_n')^2
-
2x_n^{-2}
\right)
=0 .
\]


Therefore,

\[
\boxed{
\frac12(x_n')^2
-
\frac2{x_n^2}
=
C ,
}
\]

or equivalently,

\[
\boxed{
(x_n')^2
-
\frac4{x_n^2}
=
C .
}
\]


Thus the optimal boundary-normal paths are characterized by the
same conserved energy relation appearing in Hamilton's analysis.

In particular, the absorbing boundary creates a repulsive inverse
square potential

\[
V(x_n)=\frac1{x_n^2},
\]

which prevents finite-action paths from crossing the boundary.



%%%%%%%%%%%%%%%%%%%%%%%%%%%%%%%%%%%%%%%%%%%%%%%%
\subsection{Geometric interpretation}


The above calculation shows that the Li--Yau correction term

\[
-\frac1{x_n^2}
\]

has a direct geometric meaning.

It produces an effective potential in the Harnack variational
principle:

\[
\mathcal A[x]
=
\int
\left(
\frac14|x'|^2
+
\frac1{
d(x,\partial\mathbb R^n_+)^2
}
\right)dt .
\]


Hence the boundary influence is not merely an analytic correction
term, but modifies the geometry of optimal diffusion paths.

This observation suggests a natural extension to general convex
domains, where one expects the Euclidean distance function

\[
d(x,\partial\Omega)
\]

to replace $x_n$ and the corresponding optimal paths to satisfy a
boundary-modified geodesic equation.


%%%%%%%%%%%%%%%%%%%%%%%%%%%%%%%%%%%%%%%%%%%%%%%%
\section{Sharpness and Hamilton Recovery}
%%%%%%%%%%%%%%%%%%%%%%%%%%%%%%%%%%%%%%%%%%%%%%%%


We next discuss the sharpness of the boundary correction term and
the relation with Hamilton's one-dimensional estimate.


\subsection{Sharpness of the boundary correction}



\begin{proposition}[Sharpness of the inverse-square correction]

The boundary correction term

\[
-\frac1{|x|^2}
\]

in Theorem~3.1 is of optimal order near the Dirichlet boundary.

More precisely, for the Dirichlet heat kernel itself, the quantity

\[
\Delta \log K_D(x,y,t)
+
\frac n{2t}
\]

has an inverse-square singular behavior as
$x\rightarrow \partial\mathbb R^n_+$.

\end{proposition}

\begin{proof}

Consider the Dirichlet heat kernel

\[
K_D(x,y,t)
=
\frac1{(4\pi t)^{n/2}}
\left(
e^{-\frac{|x-y|^2}{4t}}
-
e^{-\frac{|x+y|^2}{4t}}
\right).
\]


Near the boundary, the reflected Gaussian expansion gives

\[
K_D(x,y,t)
\sim
C(y,t)x_n ,
\qquad x_n\rightarrow0 .
\]

Hence

\[
\log K_D(x,y,t)
=
\log x_n+\log C(y,t)+o(1).
\]

Differentiating twice in the normal direction yields

\[
\partial_{x_nx_n}
\log K_D(x,y,t)
\sim
-\frac1{x_n^2}.
\]

Therefore the logarithmic Hessian necessarily contains an
inverse-square singularity generated by the vanishing Dirichlet
boundary condition.

Consequently, the correction term

\[
-\frac1{|x|^2}
\]

has the correct scaling order and cannot in general be removed.

\end{proof}





\begin{remark}

The inverse-square correction is consistent with the scaling
properties of the heat equation.

Under the parabolic scaling
\[
x\mapsto \lambda x,
\qquad
t\mapsto \lambda^2t ,
\]
the two quantities
\[
\frac1t
\qquad\text{and}\qquad
\frac1{|x|^2}
\]
have the same homogeneity.

Therefore the additional boundary term is the natural
scale-invariant correction associated with a codimension-one
absorbing boundary.

\end{remark}



%%%%%%%%%%%%%%%%%%%%%%%%%%%%%%%%%%%%%%%%%%%%%%%%
\subsection{Recovery of Hamilton's estimate}


We now show that the classical one-dimensional result of Hamilton is
contained in our theorem.


\begin{corollary}[Hamilton recovery]

Let $w(x,t)$ be a positive solution of the Dirichlet heat equation
on the half-line

\[
(0,\infty)\times(0,\infty),
\]

namely,

\[
w_t=w_{xx},
\qquad
w(0,t)=0 .
\]


Then

\[
(\log w)_{xx}
+
\frac1{2t}
+
\frac1{x^2}
\geq0 .
\]

Equivalently,

\[
\boxed{
l_{xx}
+
\frac1{2t}
+
\frac1{x^2}
\geq0 ,
}
\]

where

\[
l=\log w .
\]

This coincides with the simplified Li--Yau estimate obtained by
Hamilton.

\end{corollary}



\begin{proof}


When

\[
n=1,
\]

the half-space becomes the half-line

\[
\mathbb R_+=(0,\infty).
\]

Theorem~3.1 gives

\[
\Delta\log w
\geq
-\frac1{2t}
-\frac1{x^2}.
\]

Since in one dimension

\[
\Delta=\partial_{xx},
\]

we obtain

\[
(\log w)_{xx}
\geq
-\frac1{2t}
-\frac1{x^2}.
\]

Rearranging gives

\[
(\log w)_{xx}
+
\frac1{2t}
+
\frac1{x^2}
\geq0 .
\]

This is precisely Hamilton's simplified boundary Li--Yau
estimate.

\end{proof}

\begin{thebibliography}{99}
\bibitem[Hamilton2011]{Hamilton2011} R. S. Hamilton, \textit{Li--Yau Estimates and Their Harnack Inequalities}, in \textit{Geometry and Analysis}, International Press, 2011, 329--362.
\end{thebibliography}
\end{document}
